\documentclass[11pt]{amsart}

\usepackage[T1]{fontenc}
\usepackage{lmodern}
\usepackage{microtype}
\usepackage{mathtools,amssymb}
\usepackage[hidelinks]{hyperref}

\hypersetup{
  pdftitle={Conjecture 4.2 of Altinisik and Altintas Does Not Extend to Ten Elements},
  pdfsubject={Singular LCM matrices on gcd-closed sets},
  pdfkeywords={LCM matrix, gcd-closed set, Boolean lattice, meet semilattice}
}

\newtheorem{theorem}{Theorem}
\newtheorem{proposition}[theorem]{Proposition}
\newtheorem{lemma}[theorem]{Lemma}

\title[Conjecture 4.2 does not extend to ten elements]
{Conjecture~4.2 of Alt\i n\i\c{s}\i k and Alt\i nta\c{s}
Does Not Extend to Ten Elements}

\date{}

\subjclass[2020]{Primary 11C20; Secondary 15B36, 06A12}
\keywords{LCM matrix, gcd-closed set, Boolean lattice, meet semilattice}

\begin{document}

\begin{abstract}
Alt\i n\i\c{s}\i k and Alt\i nta\c{s} conjectured that a singular
LCM matrix on a nine-element gcd-closed set forces a cube sublattice.
Mattila, Haukkanen, and M\"antysalo proved the nine-element statement
and left open whether its conclusion extends to cardinalities
$10,11,12$. We give a ten-element gcd-closed set whose LCM matrix is
singular although its divisibility poset contains no copy of the
Boolean lattice $B_3$. An elementary stabilization yields such sets
in every cardinality $n\ge 10$, and known results for $n\le 9$ show
that ten is sharp.
\end{abstract}

\maketitle

\section{Introduction}

For a finite set $S=\{x_1,\dots,x_n\}$ of distinct positive integers,
write
\[
 [S]=\bigl(\operatorname{lcm}(x_i,x_j)\bigr)_{i,j=1}^n.
\]
The set $S$ is \emph{gcd-closed} if
$\gcd(x_i,x_j)\in S$ for all $i,j$. We call $(S,\mid)$
\emph{$B_3$-free} if no subset of $S$, with its inherited divisibility
order, is isomorphic to the eight-element Boolean lattice $B_3$.
In particular, a $B_3$-free poset contains no cube subsemilattice.

Conjecture~4.2 of Alt\i n\i\c{s}\i k and Alt\i nta\c{s}
\cite{AltinisikAltintas2017} states that if $|S|=9$ and $[S]$ is
singular, then $(S,\mid)$ contains a cube sublattice. Mattila,
Haukkanen, and M\"antysalo proved the nine-element statement, produced
a thirteen-element counterexample to its cardinality extension, and
explicitly left the cases $10,11,12$ open
\cite{MattilaHaukkanenMantysalo2024}. The result below settles those
cases. It is a counterexample to the cardinality extension of
Conjecture~4.2, not to its literal nine-element formulation.

\begin{theorem}\label{thm:main}
For an integer $n\ge 1$, there exists an $n$-element gcd-closed set
$S$ such that $[S]$ is singular and $(S,\mid)$ is $B_3$-free if and
only if $n\ge 10$.
\end{theorem}

\section{The counterexample}

Let
\[
 S_0=\{1,2,3,5,7,14,15,35,102,3570\},
 \qquad T=3570.
\]
Set
\[
 A=\{2,3,5,7\},\qquad M=\{14,15,35,102\}.
\]
The elements of $M$ lie above the atom pairs
\[
 \{2,7\},\quad \{3,5\},\quad \{5,7\},\quad \{2,3\},
\]
respectively. Their atom-pair graph is the cycle
$2-3-5-7-2$.

\begin{proposition}\label{prop:base}
The set $S_0$ is gcd-closed, the matrix $[S_0]$ is singular, and
$(S_0,\mid)$ is $B_3$-free.
\end{proposition}

\begin{proof}
Since
\[
 102=2\cdot3\cdot17,
 \qquad T=2\cdot3\cdot5\cdot7\cdot17,
\]
every element of $S_0$ divides $T$. Distinct elements of $A$ have
gcd $1$, and for $p\in A$ and $m\in M$ one has
$\gcd(p,m)=p$ if $p\mid m$ and $1$ otherwise. Among the elements of
$M$,
\[
\begin{aligned}
 \gcd(14,15)&=1, & \gcd(14,35)&=7, & \gcd(14,102)&=2,\\
 \gcd(15,35)&=5, & \gcd(15,102)&=3, & \gcd(35,102)&=1.
\end{aligned}
\]
Thus $S_0$ is gcd-closed.

In the displayed order of $S_0$, let
\[
 v=(-3570,1785,1190,714,510,-255,-238,-102,-35,1)^{\mathsf T}.
\]
Equivalently,
\[
 v_x=\varepsilon_x\frac{T}{x},
 \qquad
 \varepsilon_x=
 \begin{cases}
 -1,&x\in\{1\}\cup M,\\
  1,&x\in A\cup\{T\}.
 \end{cases}
\]
For $s\in S_0$,
\[
 ([S_0]v)_s
 =sT\sum_{x\in S_0}\frac{\varepsilon_x}{\gcd(s,x)}.
\]
The sum on the right is zero for every row type. Indeed,
\[
 -1+4-4+1=0
\]
when $s=1$. If $s=p\in A$, then, since every atom has degree two in
the atom-pair graph, the sum is
\[
 -1+\left(3+\frac1p\right)
 -\left(2+\frac2p\right)+\frac1p=0.
\]
If $s=m\in M$ lies above $p,q\in A$, it is
\[
 -1+\left(2+\frac1p+\frac1q\right)
 -\left(1+\frac1p+\frac1q+\frac1m\right)+\frac1m=0.
\]
Finally, for $s=T$ it is
\[
 -1+\frac12+\frac13+\frac15+\frac17
 -\frac1{14}-\frac1{15}-\frac1{35}-\frac1{102}
 +\frac1{3570}=0,
\]
as is seen after multiplication by $3570$:
\[
 -3570+4199-630+1=0.
\]
Hence $[S_0]v=0$, so $[S_0]$ is singular.

The divisibility poset has four levels,
\[
 \{1\},\qquad A,\qquad M,\qquad \{T\}.
\]
Every element of $B_3$ lies on a maximal chain of four elements,
and every four-element chain in $(S_0,\mid)$ uses one element from
each displayed level. Therefore any copy of $B_3$ in $(S_0,\mid)$
would use $1$ and $T$ as its bottom and top, three elements of $A$ as
its atoms, and three elements of $M$ as its coatoms. Those coatoms
would represent the three pairs of the selected atoms and hence
would form a triangle in the atom-pair graph. That graph is the
triangle-free cycle $C_4$. Thus $(S_0,\mid)$ is $B_3$-free.
\end{proof}

\section{Larger orders and sharpness}

The singularity-preserving construction below goes back to Hong
\cite{Hong1999}; we record that it also preserves $B_3$-freeness.

\begin{lemma}\label{lem:stabilization}
Let $S$ be a finite gcd-closed set containing $1$, and let $a>1$ be an
integer. Put
\[
 \sigma_a(S)=\{1\}\cup aS.
\]
If $[S]$ is singular, then $[\sigma_a(S)]$ is singular. If
$(S,\mid)$ is $B_3$-free, then so is $(\sigma_a(S),\mid)$.
\end{lemma}

\begin{proof}
Gcd-closure follows from $\gcd(1,ax)=1$ and
$\gcd(ax,ay)=a\gcd(x,y)$, and
$|\sigma_a(S)|=|S|+1$.
Write $S=\{x_1,\dots,x_n\}$ with $x_1=1$ and
$x=(x_1,\dots,x_n)^{\mathsf T}$. In the order
$1,ax_1,\dots,ax_n$,
\[
 [\sigma_a(S)]
 =
 \begin{pmatrix}
  1  & ax^{\mathsf T}\\
  ax & a[S]
 \end{pmatrix}.
\]
Choose $0\ne z\in\ker[S]$. The row of $[S]$ indexed by $1$ is
$x^{\mathsf T}$, so $x^{\mathsf T}z=0$. Consequently,
\[
 [\sigma_a(S)]\binom{0}{z}=0.
\]

Suppose that $Q\subseteq\sigma_a(S)$ is a copy of $B_3$. If
$1\notin Q$, then $Q\subseteq aS$, and division by $a$ gives a copy of
$B_3$ in $S$. If $1\in Q$, then $1$ is the bottom of $Q$. The element
$a$ cannot belong to $Q$: within $Q$ it would have six strict upper
bounds, whereas a nonbottom element of $B_3$ has at most three.
Thus
\[
 Q'=(Q\setminus\{1\})\cup\{a\}
\]
is a copy of $B_3$ in $aS$, and division by $a$ gives one in $S$.
Thus $B_3$-freeness is preserved.
\end{proof}

\begin{proof}[Proof of Theorem~\ref{thm:main}]
Proposition~\ref{prop:base} gives the case $n=10$, and iterating
Lemma~\ref{lem:stabilization} gives every $n\ge 10$.

For the converse, Hong proved that every LCM matrix on a gcd-closed
set with at most seven elements is nonsingular \cite{Hong1999}.
Haukkanen, Mattila, and M\"antysalo proved that an eight-element
singular example has divisibility poset isomorphic to $B_3$
\cite{HaukkanenMattilaMantysalo2015}. Mattila, Haukkanen, and
M\"antysalo proved that every nine-element singular example contains
a cube subsemilattice \cite{MattilaHaukkanenMantysalo2024}.
Consequently no $B_3$-free singular example has fewer than ten
elements.
\end{proof}

\begin{thebibliography}{9}

\bibitem{AltinisikAltintas2017}
E.~Alt\i n\i\c{s}\i k and T.~Alt\i nta\c{s},
\emph{A note on the singularity of LCM matrices on GCD-closed sets
with 9 elements},
Journal of Science and Arts, no.~3(40) (2017), 413--422.

\bibitem{HaukkanenMattilaMantysalo2015}
P.~Haukkanen, M.~Mattila, and J.~M\"antysalo,
\emph{Studying the singularity of LCM-type matrices via semilattice
structures and their M\"obius functions},
J. Combin. Theory Ser. A \textbf{135} (2015), 181--200,
\href{https://doi.org/10.1016/j.jcta.2015.05.002}
{doi:10.1016/j.jcta.2015.05.002}.

\bibitem{Hong1999}
S.~Hong,
\emph{On the Bourque--Ligh conjecture of least common multiple
matrices},
J. Algebra \textbf{218} (1999), 216--228,
\href{https://doi.org/10.1006/jabr.1998.7844}
{doi:10.1006/jabr.1998.7844}.

\bibitem{MattilaHaukkanenMantysalo2024}
M.~Mattila, P.~Haukkanen, and J.~M\"antysalo,
\emph{Some further applications of a lattice theoretic method in the
study of singular LCM matrices},
Linear Algebra Appl. \textbf{693} (2024), 98--115,
\href{https://doi.org/10.1016/j.laa.2023.04.027}
{doi:10.1016/j.laa.2023.04.027}.

\end{thebibliography}

\end{document}